% presentacion-base.tex
% Plantilla base de presentaciones — Económetra Lab, Universidad de Pamplona
%
% Cómo usarla:
%   1. Copia este archivo (y econometralab.sty) a la carpeta de la unidad que vayas a preparar,
%      p. ej. respaldo/estadistica-descriptiva/unidad1.tex
%   2. Ajusta los campos de la sección "DATOS DE LA PRESENTACIÓN" más abajo.
%   3. Reemplaza el contenido de ejemplo por el material real de la unidad.
%   4. Compila con pdflatex (dos veces si usas \tableofcontents o referencias cruzadas):
%        pdflatex presentacion-base.tex
%
% Requiere una distribución LaTeX estándar (TeX Live / MiKTeX) con beamer y tikz — no
% depende de ningún paquete o tema externo adicional.

\documentclass[aspectratio=169,11pt]{beamer}
\usepackage[spanish,es-noquoting]{babel}
\usepackage{econometralab}

% ============================================================================
% DATOS DE LA PRESENTACIÓN — edita esto para cada unidad/curso
% ============================================================================
\setElCourse{Estadística Descriptiva}
\setElUnit{Unidad 1 \textperiodcentered\ Conceptos básicos y descripción de las variables}
\setElProfessor{Sadan De la Cruz}
\setElTerm{Universidad de Pamplona \textperiodcentered\ 2026-II}
\setElCourseShort{Estadística Descriptiva}

\title{Estadística Descriptiva}
\author{Sadan De la Cruz}
\institute{Departamento de Economía, Universidad de Pamplona}
\date{2026-II}
% ============================================================================

\begin{document}

% ---- Portada ---------------------------------------------------------------
\maketitle

% ---- Agenda -----------------------------------------------------------------
\begin{frame}{Agenda}
  \tableofcontents
\end{frame}

% =============================================================================
\section{Introducción}
% =============================================================================

\begin{frame}{¿Por qué importa este tema?}
  \begin{itemize}
    \item Todo análisis económico riguroso empieza por \textbf{entender los datos} antes de modelarlos.
    \item Las variables económicas rara vez se explican solas: necesitan clasificación, resumen y representación adecuada.
    \item Esta unidad sienta las bases conceptuales que usaremos durante todo el curso.
  \end{itemize}
\end{frame}

\begin{frame}{Objetivos de la unidad}
  \begin{itemize}
    \item Clasificar variables según su naturaleza (cualitativas / cuantitativas).
    \item Distinguir entre escalas de medición: nominal, ordinal, de intervalo y de razón.
    \item Reconocer fuentes de datos económicos oficiales y su estructura.
  \end{itemize}
\end{frame}

% =============================================================================
\section{Marco conceptual}
% =============================================================================

\begin{frame}{Definición}
  \begin{block}{Variable estadística}
    Una \textbf{variable} es una característica observable que puede tomar distintos valores
    en las unidades de un conjunto de estudio (individuos, hogares, municipios, empresas, etc.).
  \end{block}
  \vspace{0.3cm}
  \begin{alertblock}{Ejemplo}
    En la Gran Encuesta Integrada de Hogares (GEIH), el \emph{nivel educativo} es una variable
    cualitativa ordinal; el \emph{ingreso laboral mensual} es una variable cuantitativa continua.
  \end{alertblock}
\end{frame}

\begin{frame}{Clasificación de variables}
  \begin{table}
    \centering
    \begin{tabular}{lll}
      \toprule
      \textbf{Tipo} & \textbf{Subtipo} & \textbf{Ejemplo} \\
      \midrule
      Cualitativa & Nominal & Sexo, sector económico \\
      Cualitativa & Ordinal & Nivel educativo, estrato \\
      Cuantitativa & Discreta & Número de hijos en el hogar \\
      Cuantitativa & Continua & Ingreso laboral, edad \\
      \bottomrule
    \end{tabular}
  \end{table}
\end{frame}

% =============================================================================
\section{Aplicación}
% =============================================================================

\begin{frame}{Caso aplicado: GEIH}
  \begin{itemize}
    \item Fuente: Gran Encuesta Integrada de Hogares, DANE.
    \item Unidad de observación: hogar / persona.
    \item Variables de interés: ingreso, ocupación, nivel educativo, municipio.
  \end{itemize}
  \vspace{0.3cm}
  \begin{block}{Pregunta para discusión}
    ¿Qué escala de medición corresponde a la variable \emph{municipio de residencia}? ¿Y a
    \emph{número de personas en el hogar}?
  \end{block}
\end{frame}

% =============================================================================
\section{Cierre}
% =============================================================================

\begin{frame}{Para llevar}
  \begin{itemize}
    \item La naturaleza de una variable determina qué herramientas estadísticas podemos usar con ella.
    \item Clasificar bien las variables es el primer paso de cualquier análisis descriptivo.
    \item En la próxima unidad aplicaremos esta clasificación a tablas de frecuencia y gráficas.
  \end{itemize}
\end{frame}

\begin{frame}{Referencias}
  \begin{itemize}
    \item Lind, D., Marchal, W. \& Wathen, S. \emph{Estadística aplicada a los negocios y la economía}.
    \item Chao, L. \emph{Estadística para ciencias administrativas}.
  \end{itemize}
\end{frame}

\end{document}
